\documentclass[11pt,a4paper]{article}

% --- Pacotes de Idioma e Codificação ---
\usepackage[utf8]{inputenc}
\usepackage[T1]{fontenc}
\usepackage[portuguese]{babel}

% --- Design e Layout ---
\usepackage[margin=2.5cm]{geometry}
\usepackage{charter} % Fonte profissional e legível
\usepackage{lastpage}
\usepackage{fancyhdr}
\usepackage[table]{xcolor}
\usepackage{tabularx}
\usepackage{booktabs}
\usepackage{xcolor}
\usepackage{enumitem}
\usepackage{amssymb}

% --- Configuração de Cores ---
\definecolor{primaryblue}{RGB}{0, 51, 102}

% --- Cabeçalho e Rodapé ---
\pagestyle{fancy}
\fancyhf{}
\fancyhead[L]{\small \textbf{Registo de Entrevista} | EcoRide Solutions}
\fancyhead[R]{\small \thepage\ de \pageref{LastPage}}

\renewcommand{\headrulewidth}{0.4pt}

% --- Ambiente Personalizado para Transcrição ---
% Facilita a distinção entre Entrevistador (E) e Entrevistado (R)
\newlist{transcription}{description}{1}
\setlist[transcription]{
	labelsep=1em,
	leftmargin=3.5em,
	style=nextline,
	font=\bfseries\color{primaryblue}
}

\newcommand{\pergunta}[1]{\item[P:] \textit{#1}}
\newcommand{\resposta}[1]{\item[R:] #1}

% --- Início do Documento ---
\begin{document}
	
	% Título Principal
	{\noindent\LARGE\bfseries\color{primaryblue} Relatório de Entrevista Individual}
	\\[0.5cm]
	
	% --- Tabela de Metadados ---
	\renewcommand{\arraystretch}{1.5}
	\begin{tabularx}{\textwidth}{@{}|l|X|l|X|@{}}
		\hline
		\rowcolor{gray!10} \textbf{ID Entrevista} & \#002 & \textbf{Data} & 26 de Fevereiro de 2026 \\ \hline
		\textbf{Hora Início} & 10:00 & \textbf{Hora Fim} & 10:30 \\ \hline
		\textbf{Local/Meio} & Sala de Reuniões A & \textbf{Estado} & Finalizada / Transcrita \\ \hline
	\end{tabularx}
	
	\vspace{0.5cm}
	
	% --- Participantes ---
	\begin{tabularx}{\textwidth}{@{}|l|X|@{}}
		\hline
		\rowcolor{gray!10} \textbf{Intervenientes} & \textbf{Cargo / Função} \\ \hline
		\textbf{Entrevistador(a)} & Eduardo Fernandes - Analista\\ \hline
		\textbf{Entrevistado(a)}  & Silvina Matagal - Assistente Administrativa \\ \hline
	\end{tabularx}
	
	\vspace{0.8cm}
	
	\hrule % Linha separadora opcional
	\vspace{0.2cm}
	
	% --- Transcrição ---
	\section*{1. Transcrição / Notas da Entrevista}
	
	\begin{transcription}
		
		\pergunta{Pode-me dizer de uma forma geral aquilo que faz no seu dia a dia?}
		\resposta{No meu dia a dia, passo grande parte do tempo a atender clientes, tanto presencialmente como por telefone ou email. Registo as entradas das trotinetes, preparo as ordens de serviço, organizo as fichas em papel, faço orçamentos, contacto clientes sempre que há novidades e faço a ponte entre eles e os mecânicos. Também trato da parte administrativa, como preparar a informação para a faturação e manter tudo organizado para que os mecânicos saibam o que têm para fazer.}

		\pergunta{No que toca a criar ordens de serviço, que tipo de informação pede, geralmente, aos clientes?}
		\resposta{Quando estou a criar uma ordem de serviço, peço sempre o nome do cliente, o contacto telefónico ou o email e, se for para faturação, o número de contribuinte. Depois peço os dados da trotinete, como a marca, o modelo e, quando o cliente sabe, o número de série. Também escrevo o problema que o cliente descreve, porque isso ajuda os mecânicos a perceberem o que procurar. Tento sempre recolher o máximo de informação possível logo à entrada.}

		\pergunta{Para além disso, não precisa de verificar documentos antigos relativos a esse cliente?}
		\resposta{Sim, muitas vezes preciso de verificar documentos antigos. Como ainda trabalhamos com fichas em papel, tenho de ir procurar nas pastas se o cliente já veio cá antes ou se aquela trotinete já teve reparações anteriores. Isso ajuda a perceber se há algum histórico de problemas ou peças que já foram trocadas. Às vezes demora um pouco porque não está tudo centralizado.}

		\pergunta{Quando precisa de alterar um orçamento, como é que o faz? Cria uma ordem de serviço nova ou altera a que já tinha?}
		\resposta{Quando é preciso alterar um orçamento, não criamos uma ordem nova. Pegamos na mesma ficha e atualizamos o valor, normalmente riscando o anterior e escrevendo o novo, com uma nota a explicar o motivo. É tudo feito manualmente. Depois disso, tenho de contactar o cliente para pedir autorização antes de o mecânico continuar.}

		\pergunta{Como contacta o cliente para o informar sobre a alteração do orçamento?}
		\resposta{Para contactar o cliente sobre alterações no orçamento, uso quase sempre o telefone ou o email. Explico o que o mecânico encontrou e pergunto se o cliente autoriza avançar com o novo valor.}
		
		\pergunta{E se o cliente ligar a perguntar pelo estado da reparação, onde é que consulta essa informação atualmente? É fácil de encontrar?}
		\resposta{Quando um cliente liga a perguntar pelo estado da reparação, vou procurar no caderno ou nas folhas onde temos as ordens de serviço. Às vezes tenho de ir perguntar ao mecânico ou ao gerente, porque nem sempre está tudo escrito. Não é muito prático, principalmente se estiver sozinha na receção e com mais clientes à espera. Se o sistema tivesse essa informação acessível por nome ou número da ordem, seria muito mais rápido.}
		
		\pergunta{Como é que informa os mecânicos sobre as ordens de serviço que ainda não foram atendidas/reparadas?}
		\resposta{Para informar os mecânicos sobre as ordens que ainda não foram atendidas, deixo as fichas numa pasta que está numa zona específica da bancada deles. Eles já sabem que as fichas que estão ali são as próximas a serem vistas. Quando é algo urgente, aviso verbalmente.}
		
		\pergunta{Quando uma reparação é feita, como é que o sabe e como sabe quanto cobrar ao cliente e como é que o contacta?}
		\resposta{Quando uma reparação é feita, o mecânico traz-me a ficha com as anotações do que foi feito e das peças que usou. A partir daí, consigo calcular quanto o cliente tem de pagar, somando o preço dos arranjos realizados e o preço das peças. Depois contacto o cliente, seja por telefone ou email, para avisar que já pode vir buscar a trotinete. Quando o cliente chega, passo a fatura no programa que usamos e finalizo o processo.
			
		\pergunta{Quais são as formas de pagamento que mais usa (numerário, multibanco, MBWay)?}
		\resposta{As formas de pagamento mais usadas são numerário, multibanco e MBWay. Cada cliente tem a sua preferência, mas ultimamente o MBWay tem sido muito comum, especialmente entre os mais jovens. Sim, era importante que o sistema registasse qual foi o método usado em cada venda, até para depois sabermos quanto entrou em dinheiro.}
		
		\pergunta{No momento da entrega, como é que confirma que o cliente pagou tudo o que devia? O sistema deve bloquear a entrega se houver valores pendentes?}
		\resposta{No momento da entrega, confirmo se o cliente pagou tudo o que devia consultando o talão da fatura ou o registo que temos no caderno ou no programa de faturação. Se não houver registo de pagamento, não entregamos a trotinete. Já aconteceu algumas vezes o cliente dizer que depois paga, mas isso dá problemas, por isso agora só entregamos com tudo pago.}
		
		%\pergunta{}
		%\resposta{}
}
		

	\end{transcription}

%	\vspace{0.8cm}
%	
%	% --- Observações e Próximos Passos ---
%	\section*{3. Análise e Observações Gerais}
%	\begin{itemize}[leftmargin=*]
%		\item \textbf{Tom de voz:} O entrevistado demonstrou abertura, mas alguma frustração com os sistemas de IT atuais.
%		\item \textbf{Pontos Críticos:} Falta de interoperabilidade entre sistemas.
%	\end{itemize}
%	
%	\section*{4. Próximos Passos}
%	\begin{tabularx}{\textwidth}{@{}|X|l|l|@{}}
%		\hline
%		\rowcolor{gray!10} \textbf{Ação} & \textbf{Responsável} & \textbf{Prazo} \\ \hline
%		Validar fluxo de dados com equipa de IT & [Teu Nome] & 05/03/2026 \\ \hline
%		Agendar follow-up para demonstração de solução & [Teu Nome] & 12/03/2026 \\ \hline
%	\end{tabularx}

	\hrule % Linha separadora opcional
	\vspace{0.2cm}
	
	\section*{2. Requisitos obtidos}
	
		\begin{tabularx}{\textwidth}{@{}|l|X|@{}}
			\hline
			\rowcolor{gray!10} \textbf{Número} & \textbf{Descrição} \\ \hline
			\textbf{1} &  O sistema deve permitir o registo imediato da entrada de uma trotinete, incluindo dados do cliente, dados do equipamento e descrição do problema. O processo deve ser otimizado para ser executado em segundos, reduzindo filas e eliminando erros de transcrição. \\ \hline
			\textbf{2} &  O sistema deve permitir criar ordens de serviço contendo: nome do cliente, contacto (email ou telefone), NIF (opcional), marca, modelo, número de série (quando disponível) e descrição do problema. Todos os campos devem ser estruturados e validados. \\ \hline
			\textbf{3} &  O sistema deve permitir consultar rapidamente o histórico de reparações de um cliente ou de uma trotinete, incluindo serviços anteriores, peças substituídas e problemas recorrentes. \\ \hline
			\textbf{4} &  O sistema deve permitir atualizar um orçamento já criado, mantendo histórico de alterações, motivo da alteração e valor anterior. Não deve ser criada uma nova ordem de serviço. \\ \hline
			\textbf{5} &  O sistema deve permitir registar a aprovação do cliente (via telefone, email ou outro meio), associando-a à ordem de serviço antes de permitir que o mecânico continue a reparação.  \\ \hline
			\textbf{6} & O sistema deve permitir consultar o estado da reparação por nome do cliente, número da ordem ou número de série, apresentando informação atualizada. \\ \hline
			\textbf{7} &  O sistema deve disponibilizar aos mecânicos uma lista digital das ordens pendentes, ordenadas por ordem de chegada. \\ \hline
			\textbf{8} &  O sistema deve permitir ao mecânico registar digitalmente o trabalho realizado e as peças utilizadas. \\ \hline
			\textbf{9} &   O sistema deve calcular automaticamente o valor final da reparação com base nas reparações feitas, peças utilizadas e preços definidos pelo gerente. \\ \hline
			\textbf{10} &  O sistema deve permitir contactar o cliente diretamente (telefone, email ou outro canal) para informar que a trotinete está pronta para levantamento. \\ \hline
			\textbf{11} &  O sistema deve permitir registar o método de pagamento utilizado (numerário, multibanco, MBWay), para controlo financeiro e reconciliação diária.  \\ \hline
			\textbf{12} &  O sistema deve impedir a conclusão da entrega da trotinete caso existam pagamentos pendentes, exibindo alerta claro à secretária. \\ \hline
		\end{tabularx}
	
\end{document}